\section{Method: Trajectory-Anchored Dynamic Selection (TADS)}
\label{sec:method}

We formalise the data-selection problem
(Sec.~\ref{sec:problem-setup}) and the variance issue that motivates
a structural anchor (Sec.~\ref{sec:method-motivation}), introduce
the trajectory-anchor extraction
(Sec.~\ref{sec:method}) and the selection score
(Sec.~\ref{sec:method}), describe the full algorithm
(Sec.~\ref{sec:method}) together with an implementation note
on the variance-ratio weight
(Sec.~\ref{sec:method-implementation}), and finally state the
stability theorem (Sec.~\ref{sec:method-theory}).


\subsection{Problem setup}
\label{sec:problem-setup}

Let $\mathcal{D}=\{x_i\}_{i=1}^N$ be a candidate instruction-tuning
pool, where each $x_i$ is an instruction-response pair tokenized as
a single sequence. Let $f_\theta$ be a base LLM with parameters
$\theta$. Given a selection budget $B \le N$, or equivalently a
selection ratio $\rho=B/N$ with $0<\rho\le1$, the goal is to choose
a subset $\mathcal{S}\subset\mathcal{D}$ with $|\mathcal{S}|=B$ so
that fine-tuning $f_\theta$ on $\mathcal{S}$ maximizes performance
on evaluation tasks $\{T_1,\ldots,T_M\}$. In our experiments, these
tasks span multiple capabilities, including knowledge, reasoning,
code, and multilingual understanding. We work in the \emph{multi-task}
setting, where a single selected subset $\mathcal{S}$ is used for
all $M$ tasks rather than re-selected for each task. This reflects
the practical setting in which a model is fine-tuned once and then
deployed across many downstream uses.

\paragraph{Token-level state dynamics.}
For a hidden layer $l \in \{1,\ldots,L\}$ of $f_\theta$ and an
input sequence $x$ with $K_x$ valid tokens, let
$h_l^{(k)}(x;\theta)\in\mathbb{R}^d$ denote the layer-$l$ hidden
state at the $k$-th token position, where valid tokens exclude
padding tokens and $d$ is the model's hidden dimension. A causal
language model processes $x$ through a token-by-token trajectory of
hidden states. The transition between consecutive token positions is
\begin{equation}
  s_l^{(k)}(x;\theta)
  :=
  h_l^{(k)}(x;\theta)-h_l^{(k-1)}(x;\theta),
  \qquad k=2,\ldots,K_x .
  \label{eq:transition}
\end{equation}
Recent work~\citep{tang2025sde} shows that these hidden-state
transitions can contain reasoning information that is not preserved
by surface tokens alone. For data selection, however, each candidate
must be reduced to a comparable scalar score, so using the full
token-level trajectory is neither practical nor necessary. We
therefore use two complementary sample-level summaries of this
trajectory. 

\paragraph{Contextualization delta.}
Let
$h_l^{\mathrm{first}}(x;\theta) := h_l^{(1)}(x;\theta)$ and
$h_l^{\mathrm{last}}(x;\theta) := h_l^{(K_x)}(x;\theta)$ denote the
hidden states at the first and last valid token positions,
respectively.\footnote{Across standard tokenization schemes, the
first and last valid token positions typically correspond to a
sequence-start marker and a final content or sequence-end token.
These markers themselves are not the focus; what matters is the
hidden state measured at those positions. Under causal attention,
$h_l^{\mathrm{first}}$ represents the state before the content of
$x$ has been processed, while $h_l^{\mathrm{last}}$ represents the
state after the full causal pass over $x$.}
The \emph{contextualization delta} of $x$ at layer $l$ is
\begin{equation}
  \Delta h_l(x;\theta)
  :=
  h_l^{\mathrm{last}}(x;\theta) - h_l^{\mathrm{first}}(x;\theta)
  =
  \sum_{k=2}^{K_x} s_l^{(k)}(x;\theta).
  \label{eq:delta}
\end{equation}
This is the telescoping sum of token-level state transitions
(Eq.~\ref{eq:transition}) and summarizes the cumulative hidden-state
change induced by processing the full sequence. Intuitively, it
captures what the sample changes in the model's representation,
rather than where the representation lies in absolute space.
Capability-direction methods such as NAIT~\citep{nait2026} use a
similar per-sample quantity, but extract it from in-domain seed
samples. In contrast, \textsc{TADS} collects
$\{\Delta h_l(x;\theta)\}$ directly from the candidate pool itself
at each training epoch. 


  
\paragraph{Sequence-mean activation.}
The \emph{sequence-mean activation} of $x$ at layer $l$ is
\begin{equation}
  \bar h_l(x;\theta)
  :=
  \frac{1}{K_x}\sum_{k=1}^{K_x} h_l^{(k)}(x;\theta),
  \label{eq:meanact}
\end{equation}
which averages the layer-$l$ hidden states over all valid token
positions. This provides a padding-invariant sequence-level
representation that does not depend on any single endpoint and
reduces the influence of local token-level idiosyncrasies, such as
punctuation marks or special tokens.

 \subsection{Why multi-task selection benefits from a structural anchor}
\label{sec:method-motivation}

RL-based dynamic selectors following the standard PPO actor--critic
recipe~\cite{schulman2017ppo, schulman2016gae} compose two
per-candidate signals into a single composite reward: a difficulty
term $\mathcal{L}_i$, defined as the response-token cross-entropy of
the current model on $x_i$, and an uncertainty term $\mathcal{H}_i$,
defined as the mean predictive entropy over response tokens. Here,
$i$ denotes the index of candidate $x_i \in \mathcal{D}$. The
composite reward is
\begin{equation}
  \mathcal{R}_i
  =
  w\,\mathcal{L}_i + (1-w)\,\mathcal{H}_i,
  \label{eq:composite}
\end{equation}
with a variance-ratio weight
\begin{equation}
  w
  =
  \frac{\mathrm{Var}_{x \in \mathcal{D}}(\mathcal{L})}
       {\mathrm{Var}_{x \in \mathcal{D}}(\mathcal{L})
        +
        \mathrm{Var}_{x \in \mathcal{D}}(\mathcal{H})
        +
        \varepsilon}.
  \label{eq:w}
\end{equation}
Both variances are taken over the entire candidate pool
$\mathcal{D}$ at the current epoch. We adopt this construction
unchanged as our base selector~\cite{sun2025dots, kung2023activeit};
the dataset-level definition of $w$ is discussed in
Sec.~\ref{sec:method-core}.
\paragraph{The multi-task variance trap.}
To see the issue, decompose the candidate pool by task groups as
$\mathcal{D}=\bigcup_{m=1}^M \mathcal{D}_m$, where
$\mathcal{D}_m$ contains the candidates associated with task $m$.
By the law of total variance, with $m$ distributed according to the
empirical task mix induced by $\mathcal{D}$,
\begin{equation}
  \mathrm{Var}_{x \in \mathcal{D}}(\mathcal{L})
  =
  \underbrace{
  \mathbb{E}_{m \sim p_{\mathcal{D}}}
  \!\left[
    \mathrm{Var}_{x \in \mathcal{D}_m}(\mathcal{L})
  \right]
  }_{\text{within-task}}
  +
  \underbrace{
  \mathrm{Var}_{m \sim p_{\mathcal{D}}}
  \!\left[
    \mathbb{E}_{x \in \mathcal{D}_m}(\mathcal{L})
  \right]
  }_{\text{between-task}} ,
  \label{eq:total-variance}
\end{equation}
and analogously for $\mathcal{H}$. The first term captures
variation among candidates within the same task, while the second
captures variation among task-level means. In a single-task pool,
the between-task term vanishes, so $w$ mainly reflects within-task
informativeness. In a multi-task pool, however, capabilities can
have different difficulty levels and different loss or entropy
scales. The between-task term can therefore dominate the pooled
variance, making $w$ depend strongly on the task mix in
$\mathcal{D}$ rather than only on per-sample informativeness. As a
result, $w$ may drift across runs even at fixed hyperparameters, and
the scalar reward $\mathcal{R}_i$ may rank candidates according to
task-level scale differences rather than by the specific capability
each sample exercises.



\paragraph{A bounded, pool-level signal absorbs the drift.}
The trajectory anchor $v_l^{(t)}$ is introduced in
Sec.~\ref{sec:method-core}. It is a pool-level summary, namely the
top principal direction of the contextualization deltas across a
probe subset, rather than a per-sample statistic. Two properties make
it useful for correcting the drift described above. Geometrically,
$v_l^{(t)}$ captures the capability subspace that the model is most
actively shaping at the current epoch and layer. This reflects the
parameter--data geometry, rather than the way the task mix balances
the two variance terms in Eq.~\ref{eq:total-variance}. Analytically,
its epoch-to-epoch change is uniformly controlled by the learning
rate. As shown in Theorem~\ref{thm:stability},
\[
  \|v_l^{(t+1)} - v_l^{(t)}\|
  \le
  \frac{2 C_\Sigma}{\gamma_{\min}}\,\eta_t .
\]
In contrast, the variance-ratio weight $w$ has no comparable
stability bound. The \textsc{TADS} score introduced in
Sec.~\ref{sec:method-core} therefore combines the task-mix-drifting
base signal $\mathcal{R}_i \cdot a_i$ with a
stability-controlled structural signal, using each signal at the
variance scale where it is informative. 
 
 

 

\begin{figure}[t]
\centering
\begin{tikzpicture}[scale=0.95, every node/.style={font=\small}]
  \draw[gray!30, very thin, -{Latex[length=2mm]}] (-0.1,0) -- (6.4,0);
  \draw[gray!30, very thin, -{Latex[length=2mm]}] (0,-0.1) -- (0,3.8);
  \node[gray!50, font=\scriptsize] at (6.5,-0.18) {$\mathbb{R}^d$};
  \filldraw[black] (0,0) circle (1.2pt);
  \foreach \x/\y in
    {3.8/1.2, 4.2/1.6, 3.5/1.8, 4.5/2.0, 3.0/1.3, 4.0/2.3, 3.7/1.5, 4.6/1.4, 3.3/1.6}
  {
    \draw[-{Latex[length=1.6mm]}, orange!60, thin, opacity=0.75]
      (0,0) -- (\x,\y);
  }
  \node[orange!80!black, font=\footnotesize, align=left, anchor=north west] at (3.5,0.85)
    {Probe deltas\\
     $\{\Delta h_l(x;\theta_{t-1})\}_{x\in\widetilde{\mathcal{D}}_t}$};
  \draw[-{Latex[length=2.6mm]}, blue!75!black, ultra thick]
    (0,0) -- (5.0,2.0);
  \node[blue!75!black, font=\bfseries] at (5.35,2.25) {$v_l^{(t)}$};
  \draw[-{Latex[length=2.2mm]}, blue!35, dashed, semithick]
    (0,0) -- (5.0,1.65);
  \node[blue!50, font=\scriptsize] at (5.40,1.45) {$v_l^{(t-1)}$};
  \node[gray!60!black, font=\scriptsize, align=left] at (6.4,1.85)
    {Thm.~\ref{thm:stability}:\\[-1pt]
     $\le \tfrac{2C_\Sigma}{\gamma_{\min}}\,\eta_t$};
  \filldraw[red!75!black] (2.2,2.85) circle (2pt);
  \node[red!75!black, font=\footnotesize] at (1.45,3.05)
    {$h_l^{\mathrm{last}}(x_i)$};
  \draw[red!75!black, dotted, thick] (2.2,2.85) -- (2.88,1.15);
  \filldraw[red!75!black] (2.88,1.15) circle (1.5pt);
  \node[red!75!black, font=\scriptsize, anchor=west]
    at (1.7, 1.2) {$\mathrm{align}_i^{(t)}$};
\end{tikzpicture}
\caption{\textbf{Trajectory anchor at epoch~$t$.}
Orange arrows are the contextualisation deltas
$\Delta h_l(x;\theta_{t-1})$ from a probe subset
$\widetilde{\mathcal{D}}_t$---the \emph{sequence-level} trajectory
inside each sample. The bold blue arrow $v_l^{(t)}$ is their top
principal direction (the anchor); the dashed $v_l^{(t-1)}$ is the
previous epoch's anchor, the \emph{epoch-level} trajectory whose
deviation is bounded by Theorem~\ref{thm:stability}. A candidate's
alignment is the projection of $h_l^{\mathrm{last}}(x_i)$ onto
$v_l^{(t)}$.}
\label{fig:anchor}
\end{figure}



\subsection{Trajectory-Anchored Selection}
\label{sec:method-core}

\textsc{TADS} scores each candidate at epoch~$t$ by reweighting a
base dynamic selector with a structural anchor extracted from the
model's own training trajectory. We describe anchor extraction, the
selection score, and the full algorithm in turn.

\paragraph{Probe subset and trajectory anchor.}
At the start of epoch~$t$, we draw a probe subset
$\widetilde{\mathcal{D}}_t \subset \mathcal{D}$ uniformly at random
with $|\widetilde{\mathcal{D}}_t| = 1024$. The probe subset is
re-sampled at every epoch so that the anchor is not biased by a
fixed set of candidates. For each layer $l$, we collect the
contextualization deltas
$\{\Delta h_l(x;\theta_{t-1}) : x \in \widetilde{\mathcal{D}}_t\}$
from Eq.~\ref{eq:delta} and form their empirical mean and centered
covariance:
\begin{align}
  \Delta \bar{h}_l^{(t)}
   &:=
   \frac{1}{|\widetilde{\mathcal{D}}_t|}
   \sum_{x \in \widetilde{\mathcal{D}}_t}
   \Delta h_l(x;\theta_{t-1}),
   \label{eq:mu} \\[2pt]
  \Sigma_l^{(t)}
   &:=
   \frac{1}{|\widetilde{\mathcal{D}}_t|}
   \sum_{x \in \widetilde{\mathcal{D}}_t}
   \bigl(\Delta h_l(x;\theta_{t-1}) - \Delta \bar{h}_l^{(t)}\bigr)
   \bigl(\Delta h_l(x;\theta_{t-1}) - \Delta \bar{h}_l^{(t)}\bigr)^{\!\top}.
   \label{eq:sigma}
\end{align}
The trajectory anchor for layer $l$ is the top unit eigenvector of
$\Sigma_l^{(t)}$:
\begin{equation}
  v_l^{(t)}
  =
  \arg\max_{\|v\| = 1}
  v^{\!\top}\Sigma_l^{(t)}v .
  \label{eq:anchor-def}
\end{equation}
We calibrate its sign so that
$\langle v_l^{(t)}, \Delta \bar{h}_l^{(t)} \rangle > 0$, removing
the PCA sign ambiguity between $v$ and $-v$.\footnote{This sign
calibration also serves as Assumption~3 of
Theorem~\ref{thm:stability}; see Sec.~\ref{sec:method-theory}.}
Intuitively, $v_l^{(t)}$ is the direction that best explains how
layer $l$ responds to the probe samples at the current parameter
state. It represents the capability direction that the model is most
actively shaping at this epoch. The anchor depends only on
$\theta_{t-1}$ and the candidate pool; it does not require
in-domain seed samples, labels, or an auxiliary model.


\paragraph{Candidate alignment.}
For each candidate $x_i$, we measure how strongly its sequence-level
representation aligns with the trajectory anchor. At layer $l$, this
is the projection of the sequence-mean activation
$\bar{h}_l(x_i;\theta_{t-1})$ from Eq.~\ref{eq:meanact} onto
$v_l^{(t)}$. We aggregate the alignment across layers as
\begin{equation}
  \mathrm{align}_i^{(t)}
  :=
  \frac{1}{L}\sum_{l=1}^{L}
  \bigl\langle
    \bar{h}_l(x_i;\theta_{t-1}),\, v_l^{(t)}
  \bigr\rangle .
  \label{eq:align-raw}
\end{equation}
The resulting scores are then min--max normalized across the
candidate pool:
\begin{equation}
  \widetilde{\mathrm{align}}_i^{(t)}
  :=
  \frac{
    \mathrm{align}_i^{(t)}
    -
    \min\limits_{j}\,\mathrm{align}_j^{(t)}
  }{
    \max\limits_{j}\,\mathrm{align}_j^{(t)}
    -
    \min\limits_{j}\,\mathrm{align}_j^{(t)}
  }
  \in [0,1].
  \label{eq:align}
\end{equation}
Here, $i$ and $j$ index candidates in $\mathcal{D}$. Averaging over
valid token positions makes $\bar{h}_l$ padding-invariant, and
averaging over layers gives a single alignment score for each
candidate. The choice of layer subset used for alignment is examined
in App.~\ref{app:inprogress}.

\paragraph{Selection score.}
For each candidate $x_i$ at refresh step $t$, the actor--critic
pipeline provides the composite reward $\mathcal{R}_i^{(t)}$ from
Eq.~\ref{eq:composite} and a PPO action $a_i^{(t)} \in [0,1]$.
\textsc{TADS} forms the final selection score by multiplying this
reward--actor consensus base by the anchor-based alignment factor:
\begin{equation}
\boxed{
s_i^{(t)}
=
\underbrace{\mathcal{R}_i^{(t)} \cdot a_i^{(t)}}_{\substack{\text{base factor}\\[1pt]
\text{high-variance sample-level signal}}}
\cdot
\underbrace{
\left(1+\lambda \cdot \widetilde{\mathrm{align}}_i^{(t)}\right)
}_{\substack{\text{anchor factor in }[1,1+\lambda]\\[1pt]
\text{stability-controlled pool-level signal}}}
}
\label{eq:score}
\end{equation}
where $\lambda \ge 0$ controls the strength of the anchor. The
multiplicative form carries three useful properties. \emph{(i) Clean
ablation:} setting $\lambda=0$ makes the anchor factor equal to $1$,
which recovers the fixed reward--actor consensus base ranking
$\mathcal{R}_i^{(t)}a_i^{(t)}$. Thus, any gain for $\lambda>0$ is
attributable to the trajectory anchor. \emph{(ii) Variance-scale
decoupling:} the base factor inherits the task-mix drift of
Eq.~\ref{eq:total-variance}, while the anchor factor is a bounded
pool-level PCA statistic that acts as a relative modulation rather
than an additive score---a separation that additive composition would
forfeit. \emph{(iii) Stability transfer:} the anchor factor is
bounded in $[1,1+\lambda]$, and its trajectory anchor is stable by
Theorem~\ref{thm:stability}. Therefore, the anchor-induced change in
$s_i^{(t)}$ is bounded by $O(\lambda\eta_t)$; the proof is given in
App.~\ref{app:score-stability}.


\paragraph{Implementation note: dataset-level $w$.}
The variance-ratio weight $w$ in Eq.~\ref{eq:w} is computed at the
dataset level, using all $N$ candidates. A naive mini-batch
implementation would recompute $\mathrm{Var}(\mathcal{L})$ and
$\mathrm{Var}(\mathcal{H})$ separately inside each batch. When the
batch size is $|\mathcal{B}|=1$, which is common under memory
constraints, both per-batch variances are identically zero. In that
case, $w=0$, and the composite reward silently reduces to
$\mathcal{R}_i=\mathcal{H}_i$ throughout training. \textsc{TADS}
avoids this failure mode by accumulating $\{\mathcal{L}_i\}$ and
$\{\mathcal{H}_i\}$ over all batches and computing $w$ once per
epoch over the full candidate pool. We include the per-batch variant
as a diagnostic experiment in App.~\ref{app:inprogress}.





\begin{algorithm}[t]
\caption{Trajectory-Anchored Dynamic Selection (TADS)}
\label{alg:tads}
\begin{algorithmic}[1]
\REQUIRE Pool $\mathcal{D}$, base model $f_{\theta_0}$, budget $B$,
         layers $\{1,\ldots,L\}$, weight $\lambda$, epochs $T$,
         probe size $|\widetilde{\mathcal{D}}|$
\FOR{$t = 1, \dots, T$}
  \STATE Sample probe set
         $\widetilde{\mathcal{D}}_t \subset \mathcal{D}$ uniformly at random
  \FOR{$l = 1, \dots, L$}
    \STATE Compute deltas
           $\{\Delta h_l(x; \theta_{t-1})\}_{x \in \widetilde{\mathcal{D}}_t}$,
           covariance $\Sigma_l^{(t)}$
           \hfill (Eqs.~\ref{eq:delta},~\ref{eq:sigma})
    \STATE $v_l^{(t)} \leftarrow$ top-1 eigenvector of $\Sigma_l^{(t)}$,
           sign-calibrated
  \ENDFOR
  \STATE For all $i$: aggregate
         $\widetilde{\mathrm{align}}_i^{(t)}$
         \hfill (Eq.~\ref{eq:align})
  \STATE Run PPO actor--critic to obtain
         $\{\mathcal{L}_i, \mathcal{H}_i, a_i\}_{i=1}^N$
  \STATE Compute pool-level $w$, composite reward $\mathcal{R}_i$
         \hfill (Eqs.~\ref{eq:w},~\ref{eq:composite})
  \STATE Score
         $s_i^{(t)} \leftarrow
          \mathcal{R}_i \cdot a_i \cdot
          (1 + \lambda \cdot \widetilde{\mathrm{align}}_i^{(t)})$
         \hfill (Eq.~\ref{eq:score})
  \STATE $\mathcal{S}_t \leftarrow$ top-$B$ samples by $s_i^{(t)}$;
         fine-tune $\theta_{t-1} \to \theta_t$ on $\mathcal{S}_t$
\ENDFOR
\RETURN $\theta_T$
\end{algorithmic}
\end{algorithm}

\subsection{Stability of the trajectory anchor}
\label{sec:method-theory}

The anchor is extracted at each epoch from a random 1024-sample
probe. A natural concern is whether this anchor reflects a stable
structural direction or merely the randomness of the probe. The
following result gives a conditional stability guarantee: if the
empirical covariance used to define the anchor changes smoothly and
has a nontrivial top-eigenvalue gap, then the anchor itself cannot
change rapidly across epochs. Its per-epoch movement is controlled
by the learning rate, and its cumulative movement is controlled by
the accumulated learning-rate mass.

\begin{theorem}[Anchor stability]
\label{thm:stability}
Fix a layer $l$ and write $\Sigma^{(t)} := \Sigma_l^{(t)}$. Let
$\eta_t$ denote the learning rate used at epoch $t$. Assume
\emph{(1)} bounded covariance change, meaning that for some
constant $C_\Sigma>0$,
\[
  \|\Sigma^{(t+1)}-\Sigma^{(t)}\|_F
  \le
  C_\Sigma \eta_t ;
\]
\emph{(2)} a positive eigenvalue gap,
$\gamma^{(t)} := \lambda_1^{(t)}-\lambda_2^{(t)}
\ge \gamma_{\min}>0$; and \emph{(3)} consistent sign calibration,
so that consecutive anchors are oriented in the same direction.
Then the sign-calibrated trajectory anchor satisfies
\begin{equation}
  \boxed{
  \bigl\|v_l^{(t+1)}-v_l^{(t)}\bigr\|
  \le
  \frac{2C_\Sigma}{\gamma_{\min}}\eta_t
  }
  \label{eq:stability-bound}
\end{equation}
and, by summing over epochs, for any $m\ge1$,
\begin{equation}
  \bigl\|v_l^{(t+m)}-v_l^{(t)}\bigr\|
  \le
  \frac{2C_\Sigma}{\gamma_{\min}}
  \sum_{s=t}^{t+m-1}\eta_s .
  \label{eq:cumulative-bound}
\end{equation}
Thus, over any finite training horizon, the remaining drift of the
anchor is bounded by the remaining learning-rate mass. In the
asymptotic case where the learning-rate tail
$\sum_{s\ge t}\eta_s$ tends to zero, the anchors become
$\varepsilon$-stationary.
\end{theorem}

The proof uses the Davis--Kahan $\sin\Theta$
theorem~\cite{daviskahan} to relate the change in covariance
$\Sigma^{(t+1)}-\Sigma^{(t)}$ to the change in the top eigenvector.
Sign calibration then converts this angular control into the
Euclidean bound in Eq.~\ref{eq:stability-bound}. The full proof,
including the interpretation of the assumptions and a score-transfer
corollary, is given in App.~\ref{app:proof-stability}. The corollary
applies to any bounded nonnegative base score, so the resulting
stability statement for the \textsc{TADS} criterion is not specific
to the RL+CR reward. In Sec.~\ref{sec:thm-empirical}, we empirically
check the quantities appearing in the theorem, including the
per-epoch anchor change and the top-eigenvalue gap across training.
We discuss the implications for reproducibility and runtime
diagnostics in Sec.~\ref{sec:discussion}.


 

